\section{Data Records}
\label{sec:data_records}

\subsection{Dataset Composition}

The dataset encompasses \ntotal catchments with delineated watersheds across Russia.
Hydrological observations include \ndischarge discharge gauging stations and \nwaterlevel water level gauges.
The high-quality subset comprises \nhighquality Grade~A gauges ($\geq$95\% completeness) during \years.
Overall, 86\% of discharge gauges meet the decent quality threshold (see Sect.~\ref{sec:data_sources_methods} for quality tier definitions).

Hydrological signatures ($N = \nsigngauges$ gauges; half-flow date: $N = \nhalfflowgauges$) span wide ranges: mean discharge \meandischarge (0.017--8.346~mm/day), BFI \meanbaseflowindex (0.210--0.903), FDC slope \meanfdcslope (0.160--13.633).

\subsection{Dataset Structure}

Table~\ref{tab:dataset_structure} summarizes the released files.
All NetCDF files include CF-compliant metadata; missing data encoded as NaN; quality flags in the discharge file distinguish observed (0), interpolated $\leq$6 days (1), suspect (2), and missing (3) values.

\begin{table}[htbp]
\centering
\caption{Dataset file structure. Total compressed size approximately 1.2~GB.}
\label{tab:dataset_structure}
\begin{tabular}{p{2.8cm}p{1.5cm}p{3.2cm}}
\toprule
\textbf{File} & \textbf{Format} & \textbf{Key variables} \\
\midrule
\texttt{boundaries.gpkg} & GeoPackage & \ntotal polygons, area, centroid, error \\
\texttt{discharge.nc} & NetCDF-4 & $Q$ (mm/d, m$^3$/s), flags; \ndischarge $\times$ 5844~d \\
\texttt{forcing.nc} & NetCDF-4 & $P$ (MSWEP), $T$, PET (ERA5-Land); \ndischarge $\times$ 5844~d \\
\texttt{attributes.csv} & CSV & 288 HydroATLAS attributes (\nattributes primary subset) \\
\texttt{signatures.csv} & CSV & \nsignatures hydro signatures \\
\texttt{water\_level/} & CSV & Daily water level (cm) for \nwaterlevel gauges \\
\bottomrule
\end{tabular}
\end{table}

\subsection{Usage Notes}

Known caveats:
\begin{itemize}
    \item Rating curve uncertainties are typically 10--20\% at high and low flows
    \item MSWEP precipitation may underestimate solid precipitation in high-latitude regions; ERA5-Land temperature and PET are reanalysis estimates with inherent model uncertainties
    \item Temporal coverage (\years, 16 years) is shorter than most CAMELS datasets, limiting multi-decadal variability analysis
    \item Gauge density is lower in remote northern and eastern regions; western European Russia is overrepresented
    \item The dataset contains nested catchments; users should account for this in spatial analyses to avoid double-counting
    \item Runoff ratio depends on precipitation product choice (see Sect.~\ref{sec:technical_validation})
\end{itemize}

The dataset is publicly available under \license license at \placeholder{[DOI: 10.XXXX/zenodo.XXXXXXX]}.
Processing code is available at \placeholder{[github.com/username/camels-ru]}.

\subsection{Limitations}
\label{sec:limitations}

\begin{itemize}
    \item \textbf{Short temporal coverage:} The 16-year record (\years) constrains multi-decadal trend analysis and captures only one phase of low-frequency climate variability.
    \item \textbf{Single DEM product:} Watershed boundaries derived from MERIT Hydro; manual verification mitigates DEM artifacts, but boundary accuracy in flat or karst terrain remains limited by the underlying elevation data.
    \item \textbf{Nested catchments:} Spatial analysis shows that the majority of gauges are nested within at least one larger catchment, with nesting depths up to 58 levels in major river systems (e.g., Ob basin). Users must account for nesting to avoid double-counting discharge or overstating sample independence.
    \item \textbf{Attribute temporal mismatch:} HydroATLAS attributes are based on circa-2000 global datasets; land cover and population density may have shifted during \years.
    \item \textbf{Stationarity assumption:} All hydrological signatures assume stationary catchment behavior, which may not hold for basins undergoing rapid land use change or permafrost degradation.
    \item \textbf{No independent discharge validation:} Discharge data are validated through internal QC and water balance consistency, but no comparison against external databases (e.g., GRDC) was performed due to limited overlap with publicly available records for Russia.
    \item \textbf{ERA5-Land circularity:} Water balance closure uses ERA5-Land precipitation, which also serves as forcing input, limiting the independence of this consistency check. Rating curve uncertainties (10--20\% at extremes) further contribute to water balance residuals.
    \item \textbf{Gauge network bias:} Arctic permafrost-dominated catchments and Pacific-draining basins are underrepresented relative to their geographic extent.
    \item \textbf{Water level data:} Water level records undergo basic QC (zero replacement, gap interpolation) but lack the automated A--F quality grading applied to discharge data. Users should assess per-gauge completeness before use.
    \item \textbf{Data source continuity:} The AIS GMVO platform, which served as the primary data source, ceased public operation in 2025. \dataset thus serves as a preservation archive for these discharge and water level records.
\end{itemize}

\subsection{Future Updates}
Planned dataset enhancements include:
\begin{itemize}
    \item Extension of temporal coverage using historical records from institute archives and OCR-digitized observation books (from the 1950s to present for a subset of gauges)
    \item Integration of additional discharge gauges with delineated watersheds
    \item Expansion of the hydropower/reservoir gauge network: additional reservoir gauges with delineated watersheds await meteorological forcing extraction and attribute computation
    \item Addition of remotely sensed variables (NDVI, snow cover extent) as dynamic attributes
\end{itemize}
